Machine-learning flow surrogates are almost universally scored by resampling their
predictions onto a uniform grid. We measure what it costs, then what it hides. On a
standard external-aerodynamics benchmark, kernel interpolation over seven scalars from the
case name---no network, no learning---gives $8.4\times$ lower volume-pressure error than a
Transolver trained on the same data on an area-uniform $128^2$ raster. Re-weighting those
identical predictions by the node measure gives $0.44\times$, and scoring per node on the
native cloud $0.21\times$: three measures in current use, and the ranking reverses. At
native resolution the interpolator loses on \textbf{every} channel, and inside the first
$0.005$ chord the exact least-squares bound---what \emph{any} weighting of the $800$
training fields can achieve---is $25\times$ worse on average, and worse on all $200$ cases.
That impossibility result holds only in the frame it is posed in. There a query lands
$7.8\times$ further from the training wall than from its own, with $35\%$ of the kernel
weight \emph{inside} the training airfoil; re-posed in a parameter-free wall-following
frame, nothing retuned, the same bound falls to $\mathbf{0.0044\times}$, below the
surrogate on every case. We claim the \emph{barrier}, not the gap: the deployed estimator
improves $112\times$ but still trails on the case-mean ($1.7\times$). On pressure, where no
rule was registered, there was never a barrier, and the frame worsens the bound on $185$ of
$200$ cases. The frame is undefined on $5\%$ of near-wall nodes, and the raster's own
round-trip error at the nodes exceeds the surrogate's by $413\times$.
